\documentclass[11pt,letterpaper]{article}
\usepackage{amsmath,amssymb,amsthm,graphicx}
\begin{document}
\section{Triangles and circles}
Figure~\ref{fig:sq} shows a square of side one, which we use as a unit cell.

\begin{figure}[htbp]
\centering
\includegraphics[width=3cm]{images/square.png}
\caption{A unit square.}\label{fig:sq}
\end{figure}

The diagonal of the unit square has length $\sqrt2$, and the circumscribed circle has radius $\tfrac{\sqrt2}{2}$. Its area is therefore $\pi r^2=\pi/2$.

\subsection{Law of cosines}
In a triangle with sides $a,b,c$ and angle $\gamma$ opposite $c$,
\begin{equation}
c^2=a^2+b^2-2ab\cos\gamma.\label{eq:cos}
\end{equation}
Setting $\gamma=\pi/2$ in \eqref{eq:cos} recovers Pythagoras.

\subsection{Areas}
\begin{table}[htbp]
\centering
\begin{tabular}{lcr}
\hline
Shape & Parameter & Area\\ \hline
Square & side $s$ & $s^2$\\
Disc & radius $r$ & $\pi r^2$\\
Triangle & base $b$, height $h$ & $\tfrac12 bh$\\ \hline
\end{tabular}
\caption{Some areas.}\label{tab:areas}
\end{table}
Table~\ref{tab:areas} lists three cases; see page~\pageref{tab:areas} for the table itself.
\end{document}
